% =====================================================================
%  Giới thiệu + bảng ký hiệu (P1)
% =====================================================================
\section{Giới thiệu}\label{sec:intro}

Trong bài toán phân loại nhị phân, khi hai lớp tách được bằng một siêu phẳng thì thường có vô số
siêu phẳng như vậy. Máy vector hỗ trợ (SVM) chọn siêu phẳng có \emph{lề} lớn nhất, tức là cách xa
nhất các điểm gần nó của cả hai lớp; phiên bản lề mềm cho phép một số mẫu vi phạm lề với chi phí
điều khiển bởi tham số $C$ \citep{cortes1995support}. Hai tính chất làm nên sức mạnh của SVM.
Thứ nhất, bài toán huấn luyện là quy hoạch toàn phương lồi, nên mọi cực tiểu địa phương đều là cực
tiểu toàn cục và kết quả không phụ thuộc vào khởi tạo hay thứ tự duyệt mẫu. Thứ hai, bài toán đối ngẫu chỉ truy
cập dữ liệu qua tích vô hướng; thay tích vô hướng bằng một hàm nhân xác định dương cho phép học biên
quyết định phi tuyến trong một không gian đặc trưng, có thể vô hạn chiều, mà không cần tính ánh xạ
đặc trưng một cách tường minh \citep{scholkopf2002learning}. Nghiệm cuối cùng chỉ phụ thuộc vào một
phần nhỏ các mẫu, gọi là \emph{vector hỗ trợ}.

Những tính chất này khiến SVM đặc biệt phù hợp với bài toán có \emph{dữ liệu cỡ nhỏ} (vài trăm đến
vài nghìn mẫu) và \emph{đặc trưng thiết kế thủ công} (vài chục chiều). Đây đúng là bối cảnh của ứng
dụng trong báo cáo: phân loại giống gạo Việt Nam từ đặc trưng hình thái, màu sắc và kết cấu của
từng hạt, với mục đích hỗ trợ kiểm tra tính xác thực của nhãn các giống gạo đặc sản. Ở quy mô này,
các mô hình nhiều tham số thường cần lượng dữ liệu lớn hơn nhiều so với mức có thể thu thập thủ
công. Ngược lại, SVM chỉ có hai siêu tham số chính ($C$ và tham số của hàm nhân) chọn được bằng
kiểm định chéo, và có các kết quả lý thuyết liên hệ độ thưa của nghiệm với sai số tổng quát hóa
(Mục~\ref{sec:B6}).

Mục tiêu của báo cáo là đi trọn từ lý thuyết, qua cài đặt, tới ứng dụng:
\begin{itemize}
  \item \textbf{Lý thuyết} (Mục~\ref{sec:theory}): chứng minh đầy đủ các kết quả nền tảng, gồm hình
        học của lề cực đại, bài toán đối ngẫu và điều kiện KKT của SVM lề mềm, sự tương đương với
        cực tiểu hóa mất mát hinge, định lý biểu diễn \citep{scholkopf2001generalized}, tính chất đóng
        của hàm nhân xác định dương, và cận sai số bỏ-một theo số vector hỗ trợ.
  \item \textbf{Cài đặt} (Mục~\ref{sec:impl}): cài đặt từ đầu thuật toán SMO \citep{platt1998smo}
        với chiến lược chọn cặp vi phạm KKT lớn nhất, hạ tọa độ đối ngẫu cho SVM tuyến tính
        \citep{hsieh2008dual}, hồi quy Ridge hàm nhân, mở rộng đa lớp và Platt scaling
        \citep{platt1999probabilistic}; đối chiếu kết quả với LIBSVM \citep{chang2011libsvm}.
  \item \textbf{Bài tập và thực nghiệm} (Mục~\ref{sec:exercises}): giải tay các ví dụ nhỏ, khảo sát
        ảnh hưởng của $C$ và $\gamma$ bằng kiểm định chéo lồng nhau, kiểm chứng thực nghiệm cận sai số
        bỏ-một, và so sánh thời gian huấn luyện của SMO với hạ tọa độ đối ngẫu.
  \item \textbf{Ứng dụng} (Mục~\ref{sec:app}): thu thập ảnh hạt gạo có kiểm soát, xử lý ảnh, trích
        chọn đặc trưng, so sánh SVM với k-NN và hồi quy logistic, đánh giá theo phiên chụp để tránh
        rò rỉ dữ liệu giữa tập huấn luyện và tập kiểm tra.
\end{itemize}
Mục~\ref{sec:conclusion} tóm tắt kết quả và hướng phát triển; Phụ lục~\ref{sec:contrib} trình bày
phân công công việc. Mã nguồn (thư viện \texttt{svmlib/}) và các script tái lập thí nghiệm
(\texttt{experiments/}) nằm trong kho mã của nhóm.

\subsection{Ký hiệu}\label{sec:notation}

\begin{table}[H]
\centering
\caption{Ký hiệu dùng trong báo cáo}\label{tab:notation}
\begin{tabular}{@{}ll@{}}
\toprule
Ký hiệu & Ý nghĩa \\
\midrule
$n$, $d$ & số mẫu huấn luyện, số chiều đặc trưng \\
$\vx_i \in \R^d$, $y_i \in \{-1, +1\}$ & mẫu thứ $i$ và nhãn của nó \\
$\vw$, $b$ & vector pháp tuyến và hệ số chặn của siêu phẳng \\
$f(\vx) = \vw\T\vx + b$ & hàm quyết định; dự đoán $\sign f(\vx)$ \\
$\xi_i \ge 0$ & biến bù (slack) của mẫu $i$ trong SVM lề mềm \\
$C > 0$ & tham số phạt của SVM lề mềm \\
$\lambda > 0$ & hệ số chính quy hóa, $C = 1/(\lambda n)$ \\
$\valpha \in \R^n$ & nhân tử Lagrange (biến đối ngẫu) \\
$k(\vx, \vz)$ & hàm nhân xác định dương \\
$\mK$, $K_{ij} = k(\vx_i, \vx_j)$ & ma trận Gram \\
$\mQ$, $Q_{ij} = y_i y_j K_{ij}$ & ma trận của bài toán đối ngẫu \\
$\Hk$, $\Hnorm{\cdot}$ & RKHS sinh bởi $k$ và chuẩn trên đó \\
$\gamma$ & tham số độ rộng của hàm nhân Gauss $\exp(-\gamma\norm{\vx - \vz}^2)$ \\
$\SV = \{i : \alpha_i > 0\}$, $\nsv = \abs{\SV}$ & tập và số vector hỗ trợ \\
$\Iup$, $\Ilow$ & tập chỉ số dùng trong điều kiện KKT của SMO \\
\bottomrule
\end{tabular}
\end{table}
